\documentclass[a4paper,11pt]{article}

\usepackage[margin=1.2cm,top=1.1cm,bottom=1.1cm]{geometry}
\usepackage[T1]{fontenc}
\usepackage{lmodern}
\usepackage{tabularx}
\usepackage{enumitem}
\usepackage{xcolor}
\usepackage{titlesec}
\usepackage{fancyhdr}
\usepackage{hyperref}
\usepackage{microtype}
\usepackage{graphicx}
\usepackage{needspace}

\definecolor{accent}{HTML}{17365D}
\definecolor{muted}{HTML}{4F5965}
\hypersetup{
  colorlinks=true,
  urlcolor=accent,
  linkcolor=accent,
  pdftitle={Zhongli Wang - Quantitative CV},
  pdfauthor={Zhongli Wang}
}

\pagestyle{fancy}
\fancyhf{}
\renewcommand{\headrulewidth}{0pt}
\renewcommand{\footrulewidth}{0pt}
\fancyfoot[C]{\small\textcolor{muted}{\thepage}}

\titleformat{\section}
  {\large\bfseries\color{accent}}
  {}{0pt}{}
  [\vspace{1pt}\color{accent}\titlerule]
\titlespacing*{\section}{0pt}{4pt}{2pt}

\setlength{\parindent}{0pt}
\setlength{\parskip}{1pt}
\setlist[itemize]{leftmargin=1.25em,nosep,topsep=0.5pt,itemsep=0.5pt}
\newcommand{\role}[4]{%
  \textbf{#1} \hfill \textbf{#2}\\
  \textit{#3} \hfill \textit{#4}\par
}
\newcommand{\project}[2]{%
  \textbf{#1}\\[-1pt]
  \textcolor{muted}{\small #2}\par
}

\begin{document}

\begin{center}
  {\LARGE\textbf{\textcolor{accent}{Zhongli Wang}}}\\[3pt]
  {\large Quantitative Research \textbar{} Machine Learning \textbar{} Financial Data \textbar{} AI}\\[3pt]
  Bielefeld, Germany \textbar{} PhD successfully defended on 21 September 2026 \textbar{} Hong Kong TTPS Approved \\[2pt]
  \href{tel:+4916095416015}{+49 160 95416015} \textbar{}
  \href{mailto:nkwangzhongli@gmail.com}{nkwangzhongli@gmail.com} \textbar{}
  \href{https://sites.google.com/view/zhongli-wang/}{sites.google.com/view/zhongli-wang}
\end{center}

\section{Profile}
PhD in Economics with MA in Banking \& Finance and a Data Science minor, plus BSc in Finance. Quantitative researcher combining Python, machine learning, agent-based modelling, and financial-data analysis to study complex and changing markets. Experience taking research questions from model design and large-scale data analysis through validation, interpretation, and communication, with published work in financial fraud detection and AI-enabled research workflows. Brings strong statistical reasoning, financial domain knowledge, and collaborative communication to AI-driven electronic trading.

\section{Core Competencies}
\begin{itemize}
  \item \textbf{Quantitative research:} Statistical learning, supervised learning, agent-based modelling, numerical methods, model calibration, sensitivity analysis
  \item \textbf{Python and ML:} Python (pandas, NumPy, scikit-learn, XGBoost), R, SQL; reproducible analysis, feature construction, model comparison, data visualization
  \item \textbf{Financial data and risk:} Financial statement analysis, fraud detection, accounting metrics, banking operations, monetary policy, risk assessment
  \item \textbf{Research ownership:} Research design, data preparation, computational implementation, validation, interpretation, and communication of evidence-based findings
  \item \textbf{AI and collaboration:} LLM and multi-agent workflows, technical documentation, presentations, cross-functional collaboration with economists and software engineers
\end{itemize}

\section{Research and Professional Experience}
\role{Doctoral Researcher, Economics}{Oct. 2021--Oct. 2026}{Bielefeld University and Universit\`a Cattolica del Sacro Cuore}{Bielefeld / Milan}
\begin{itemize}
  \item Design and implement computational research in Python, combining agent-based models, machine learning, numerical methods, and financial data to study market dynamics and firm behavior.
  \item Analyse AT\&T and T-Mobile financial data, construct model inputs, and compare alternative estimation and validation approaches for research on capacity, competition, and data markets.
  \item Conduct sensitivity analysis and robustness checks, document assumptions, and interpret results with attention to model limitations, overfitting risk, and changing economic conditions.
  \item Own projects from research question and data preparation through computational implementation, validation, written analysis, and presentations to international academic and business audiences.
\end{itemize}

\role{Research Assistant, Data Science Project}{Aug. 2022--Sept. 2022}{LucidMinds, DataVille 3000 project}{Amsterdam}
\begin{itemize}
  \item Gathered requirements and developed Python-based computational tools to model consumer behaviour and dynamic market interactions.
  \item Analysed model outputs, communicated technical findings to stakeholders, and iterated the approach using feedback and performance metrics.
\end{itemize}

\role{Helping Assistant}{June 2019--June 2021}{Quantitative Business Administration Team, University of Zurich}{Zurich}
\begin{itemize}
  \item Assisted with quantitative modelling and analysis using \textit{Mathematica}; developed teaching content for Applied Microeconomics, Corporate Finance, and Business Administration.
\end{itemize}

\role{HR Intern}{Sept. 2017--Dec. 2017}{Continental Automotive Systems, Tianjin}{Tianjin, China}
\begin{itemize}
  \item Supported recruitment initiatives by assisting in student hiring programs and coordinating interview logistics.
  \item Translated internal HR materials across English, German and Chinese, facilitating multilingual communication.
  \item Managed administrative tasks related to reimbursement processes, ensuring accuracy and compliance with internal procedures.
\end{itemize}

\needspace{7\baselineskip}
\role{Finance Internships}{2016--2017}{People's Bank of China; Agricultural Bank of China; Gansu Provincial Highway Aviation Tourism Financial Services Co.}{Lanzhou, China}
\begin{itemize}
  \item Analysed financial status reports and supported monetary policy analysis at People's Bank of China, gaining exposure to interest-rate management and reserve requirements.
  \item Conducted financial reporting analysis and process-improvement work at Agricultural Bank of China; reviewed banking operations and FMIS systems.
  \item Reviewed and summarized quartet agreements to support approval processes for monthly loan items and assisted in account reconciliation.
  \item Collaborated with department managers to analyze and interpret cases reported in annual financial statements.
\end{itemize}

 


\section{Education}
\textbf{PhD in Economics, Bielefeld University \& Universit\`a Cattolica} \hfill 2021--2026\\
EPOC Marie Sk\l odowska-Curie program.

\textbf{MA in Business and Economics, University of Zurich} \hfill 2018--2021\\
Major: Banking \& Finance; Minor: Data Science; magna cum laude.

\textbf{BSc in Finance, Nankai University} \hfill 2015--2018\\
GPA: 4.39/5.0.

\textbf{BA in German Language and Literature, Nankai University} \hfill 2014--2018\\
GPA: 5.0/5.0.

% \newpage
\section{Key Projects}
\project{Financial Fraud Detection Analysis}{Published in the \textit{International Review of Financial Analysis}, 2025}
\begin{itemize}
  \item Conducted financial-data analysis on accounting metrics and market indicators to identify fraud patterns and risk factors.
  \item Developed and compared supervised-learning fraud-detection models using XGBoost, Random Forest, and neural networks for risk-management applications.
  \item Interpreted model outputs and translated empirical evidence into recommendations for fraud prevention and risk assessment.
\end{itemize}


\project{Agentic Workflows for Economic Research}{Published on arXiv, 2025}
\begin{itemize}
  \item Designed and implemented AI-powered agentic workflows using large language models (LLMs) and multi-agent systems to automate economic research tasks.
  \item Built an end-to-end research-automation workflow and coordinated with software engineers and economists to integrate AI agents into research processes.
\end{itemize}

\section{Technical Skills}
\begin{tabularx}{\linewidth}{@{}>{\bfseries}p{3.2cm}X@{}}
Data and ML & Python (pandas, NumPy, matplotlib, scikit-learn, XGBoost), R (ggplot2, dplyr), SQL; feature construction, statistical analysis\\
Quantitative Methods & Agent-based modelling, numerical methods, statistical modelling, sensitivity analysis, model calibration, SMM, MCMC\\
Financial Research & Financial statement analysis, fraud detection, accounting metrics, banking operations, monetary policy, risk assessment\\
AI Workflows & LLMs, Microsoft AutoGen, multi-agent systems, research automation, technical documentation\\
Coding & Git, GitHub, Linux, Jupyter Notebook, LaTeX, VS Code\\
Collaboration & Cross-functional teamwork, stakeholder communication, requirements gathering, research presentations\\
\end{tabularx}

\section{Publications}
\begin{itemize}
  \item Zhao et al. (2025). ``Polytope Fraud Theory.'' \textit{Int. Review of Financial Analysis}, 97, 103734.
  \item Dawid et al. (2025). ``Agentic Workflows for Economic Research.'' arXiv preprint.
  \item \textbf{Wang, Z.} (2025). ``Heterogeneous Consumers and the Data Value Chain.'' SSRN Working Paper.
\end{itemize}

\enlargethispage{3\baselineskip}
\section{Awards \& Languages}
\textbf{Awards:} Doctoral Completion Scholarship, Bielefeld Young Researchers' Fund (2025); Outstanding Graduate, Nankai University (2018).

\textbf{Languages:} Mandarin Chinese (native), English (C1), German (C1), French (A2), Cantonese (A2).

\end{document}
